\documentclass[10pt,a4paper]{amsart}
\usepackage{amsmath,amssymb,amsthm,graphicx}
\usepackage{verbatim}
\newenvironment{reference}{\comment}{\endcomment}
%\newenvironment{reference}{}{}
\newenvironment{slogan}{\comment}{\endcomment}
\newenvironment{history}{\comment}{\endcomment}

% For commutative diagrams we use Xy-pic
\usepackage[all]{xy}

% We use 2cell for 2-commutative diagrams.
\xyoption{2cell}
% Theorem environments.
%
\theoremstyle{plain}
\newtheorem{theorem}[subsection]{Theorem}
\newtheorem{proposition}[subsection]{Proposition}
\newtheorem{lemma}[subsection]{Lemma}

\theoremstyle{definition}
\newtheorem{definition}[subsection]{Definition}
\newtheorem{example}[subsection]{Example}
\newtheorem{exercise}[subsection]{Exercise}
\newtheorem{situation}[subsection]{Situation}

\theoremstyle{remark}
\newtheorem{remark}[subsection]{Remark}
\newtheorem{remarks}[subsection]{Remarks}

\numberwithin{equation}{subsection}

% Macros
%
\def\lim{\mathop{\mathrm{lim}}\nolimits}
\def\colim{\mathop{\mathrm{colim}}\nolimits}
\def\Spec{\mathop{\mathrm{Spec}}}
\def\Hom{\mathop{\mathrm{Hom}}\nolimits}
\def\Ext{\mathop{\mathrm{Ext}}\nolimits}
\def\SheafHom{\mathop{\mathcal{H}\!\mathit{om}}\nolimits}
\def\SheafExt{\mathop{\mathcal{E}\!\mathit{xt}}\nolimits}
\def\Sch{\mathit{Sch}}
\def\Mor{\mathop{\mathrm{Mor}}\nolimits}
\def\Ob{\mathop{\mathrm{Ob}}\nolimits}
\def\Sh{\mathop{\mathit{Sh}}\nolimits}
\def\NL{\mathop{N\!L}\nolimits}
\def\CH{\mathop{\mathrm{CH}}\nolimits}
\def\proetale{{pro\text{-}\acute{e}tale}}
\def\etale{{\acute{e}tale}}
\def\QCoh{\mathit{QCoh}}
\def\Ker{\mathop{\mathrm{Ker}}}
\def\Im{\mathop{\mathrm{Im}}}
\def\Coker{\mathop{\mathrm{Coker}}}
\def\Coim{\mathop{\mathrm{Coim}}}

% Boxtimes
%
\DeclareMathSymbol{\boxtimes}{\mathbin}{AMSa}{"02}

%
% Macros for moduli stacks/spaces
%
\def\QCohstack{\mathcal{QC}\!\mathit{oh}}
\def\Cohstack{\mathcal{C}\!\mathit{oh}}
\def\Spacesstack{\mathcal{S}\!\mathit{paces}}
\def\Quotfunctor{\mathrm{Quot}}
\def\Hilbfunctor{\mathrm{Hilb}}
\def\Curvesstack{\mathcal{C}\!\mathit{urves}}
\def\Polarizedstack{\mathcal{P}\!\mathit{olarized}}
\def\Complexesstack{\mathcal{C}\!\mathit{omplexes}}
% \Pic is the operator that assigns to X its picard group, usage \Pic(X)
% \Picardstack_{X/B} denotes the Picard stack of X over B
% \Picardfunctor_{X/B} denotes the Picard functor of X over B
\def\Pic{\mathop{\mathrm{Pic}}\nolimits}
\def\Picardstack{\mathcal{P}\!\mathit{ic}}
\def\Picardfunctor{\mathrm{Pic}}
\def\Deformationcategory{\mathcal{D}\!\mathit{ef}}
\usepackage{cleveref}
\crefname{originalreference}{original reference}{original references}
\Crefname{originalreference}{Original reference}{Original references}
\makeatletter
\newlabel{situation-shriek}{{0F42}{source}{}{sourcepointer-1}{}}
\newlabel{situation-shriek@cref}{{[originalreference][0][]0F42}{[1][0][]source}{}{}{}}
\makeatother
\crefname{originalequation}{original equation}{original equations}
\Crefname{originalequation}{Original equation}{Original equations}
\makeatletter
\newlabel{equation-compare-with-pullback}{{0A9S}{source}{}{sourcepointer-2}{}}
\newlabel{equation-compare-with-pullback@cref}{{[originalequation][0][]0A9S}{[1][0][]source}{}{}{}}
\makeatother
\crefname{originallemma}{original lemma}{original lemmas}
\Crefname{originallemma}{Original lemma}{Original lemmas}
\makeatletter
\newlabel{lemma-transitivity-compare-with-pullback}{{0B6R}{source}{}{sourcepointer-3}{}}
\newlabel{lemma-transitivity-compare-with-pullback@cref}{{[originallemma][0][]0B6R}{[1][0][]source}{}{}{}}
\makeatother
\crefname{originallemma}{original lemma}{original lemmas}
\Crefname{originallemma}{Original lemma}{Original lemmas}
\makeatletter
\newlabel{lemma-shriek-well-defined}{{0AA0}{source}{}{sourcepointer-4}{}}
\newlabel{lemma-shriek-well-defined@cref}{{[originallemma][0][]0AA0}{[1][0][]source}{}{}{}}
\makeatother
\crefname{originallemma}{original lemma}{original lemmas}
\Crefname{originallemma}{Original lemma}{Original lemmas}
\makeatletter
\newlabel{lemma-twisted-inverse-image}{{0A9E}{source}{}{sourcepointer-5}{}}
\newlabel{lemma-twisted-inverse-image@cref}{{[originallemma][0][]0A9E}{[1][0][]source}{}{}{}}
\makeatother
\crefname{originallemma}{original lemma}{original lemmas}
\Crefname{originallemma}{Original lemma}{Original lemmas}
\makeatletter
\newlabel{lemma-compare-on-open}{{0B6Q}{source}{}{sourcepointer-6}{}}
\newlabel{lemma-compare-on-open@cref}{{[originallemma][0][]0B6Q}{[1][0][]source}{}{}{}}
\makeatother
\crefname{originalequation}{original equation}{original equations}
\Crefname{originalequation}{Original equation}{Original equations}
\makeatletter
\newlabel{equation-sheafy}{{0A9L}{source}{}{sourcepointer-7}{}}
\newlabel{equation-sheafy@cref}{{[originalequation][0][]0A9L}{[1][0][]source}{}{}{}}
\makeatother
\crefname{originallemma}{original lemma}{original lemmas}
\Crefname{originallemma}{Original lemma}{Original lemmas}
\makeatletter
\newlabel{lemma-restriction-compare-with-pullback}{{0B6P}{source}{}{sourcepointer-8}{}}
\newlabel{lemma-restriction-compare-with-pullback@cref}{{[originallemma][0][]0B6P}{[1][0][]source}{}{}{}}
\makeatother
\crefname{originallemma}{original lemma}{original lemmas}
\Crefname{originallemma}{Original lemma}{Original lemmas}
\makeatletter
\newlabel{lemma-upper-shriek-composition}{{0ATX}{source}{}{sourcepointer-9}{}}
\newlabel{lemma-upper-shriek-composition@cref}{{[originallemma][0][]0ATX}{[1][0][]source}{}{}{}}
\makeatother
\title{The comparison map from the pulled back object tensored with the shriek of the structure sheaf to the shriek is well defined and associative}
\author{The Stacks Project Authors}
\date{stacks-project:duality.tex:local-snapshot-20261012; GNU FDL 1.2 or later}
\begin{document}
\maketitle
\noindent\textit{Editorial note.} All delivered mathematical text is the author's own, copied from the pinned Stacks Project chapter file duality.tex; no statement or proof step was rewritten, reordered or removed.\par\medskip
\noindent\textit{Editorial note.} Preamble provenance: the pinned chapter file opens with an input of the shared preamble and carries no preamble text of its own, so the project's single Stacks-wide preamble staged at sources/stacks-duality/preamble.tex is excerpted verbatim under the recipe fixed by LANE\_BRIEF 22.1, as three spans, lines 8-13 (verbatim package and the reference/slogan/history environments), lines 14-18 (Xy-pic setup) and lines 154-217 (the native theorem declarations and mathematical macros). Nothing was added, removed or redefined; the wrapper class and the amsmath/amssymb/amsthm/graphicx base packages are supplied by the assembler.\par\medskip
\noindent\textit{Editorial note.} The counted statement and proof contain no cross reference, no citation and no included file, so the delivered body is self contained.\par\medskip
\noindent\textit{Editorial note.} Canonical location of the counted proposition: Stacks tag 0B6T (sources/stacks-duality/pinned-original-source.tex line 4427).\par\medskip
\noindent\textit{Allowed standard prerequisites (editorial).}\begin{itemize}
\item The upper shriek functor through a compactification and its well definedness. Same chapter, Lemma labelled lemma-shriek-well-defined, cited at proof lines 4470-4474 and repeatedly as the model argument, and Lemma labelled lemma-twisted-inverse-image for the right adjoint. Supplies the description f\textasciicircum{}! = j\textasciicircum{}* o a-bar and the independence argument to be imitated.
\item The comparison map between pullback and the right adjoint for a proper morphism, and its transitivity. Same chapter, equation labelled equation-compare-with-pullback, cited at proof lines 4458-4461 and 4469-4472, and Lemma labelled lemma-transitivity-compare-with-pullback, cited at proof lines 4461 and 4520-4524 and 4552-4556. Produces the maps whose restrictions define mu and proves the commutativity diagrams.
\item Base change maps for the right adjoints and their restriction to the open. Same chapter, equation labelled equation-sheafy, cited at proof lines 4516-4520, Lemma labelled lemma-compare-on-open, cited at proof lines 4526-4530, and Lemma labelled lemma-restriction-compare-with-pullback, cited at proof lines 4534-4538. Shows that sigma and tau restrict to isomorphisms and identifies the restrictions with the identity.
\item Composition of the shriek functors. Same chapter, Lemma labelled lemma-upper-shriek-composition, cited at proof lines 4544-4550. Supplies the construction of (g o f)\textasciicircum{}! -> f\textasciicircum{}! o g\textasciicircum{}! used for the hexagon.
\end{itemize}\medskip
\ifdefined\hypertarget\hypertarget{sourcepointer-1}{}\fi
\noindent\textit{Source reference.} Original reference 0F42: sources/stacks-duality/pinned-original-source.tex, source lines 3953--3953. Same-chapter pointer into the author environment situation at source line 3953; invoked at proof line None of the counted proof. It is retained only as a navigation pointer; every load-bearing step is inside the delivered span.\par\medskip
\ifdefined\hypertarget\hypertarget{sourcepointer-2}{}\fi
\noindent\textit{Source reference.} Original equation 0A9S: sources/stacks-duality/pinned-original-source.tex, source lines 2079--2079. Same-chapter pointer into the author environment equation at source line 2079; invoked at proof line 4458 of the counted proof. It is retained only as a navigation pointer; every load-bearing step is inside the delivered span.\par\medskip
\ifdefined\hypertarget\hypertarget{sourcepointer-3}{}\fi
\noindent\textit{Source reference.} Original lemma 0B6R: sources/stacks-duality/pinned-original-source.tex, source lines 2189--2189. Same-chapter pointer into the author environment lemma at source line 2189; invoked at proof line 4459 of the counted proof. It is retained only as a navigation pointer; every load-bearing step is inside the delivered span.\par\medskip
\ifdefined\hypertarget\hypertarget{sourcepointer-4}{}\fi
\noindent\textit{Source reference.} Original lemma 0AA0: sources/stacks-duality/pinned-original-source.tex, source lines 3986--3986. Same-chapter pointer into the author environment lemma at source line 3986; invoked at proof line 4474 of the counted proof. It is retained only as a navigation pointer; every load-bearing step is inside the delivered span.\par\medskip
\ifdefined\hypertarget\hypertarget{sourcepointer-5}{}\fi
\noindent\textit{Source reference.} Original lemma 0A9E: sources/stacks-duality/pinned-original-source.tex, source lines 653--653. Same-chapter pointer into the author environment lemma at source line 653; invoked at proof line 4481 of the counted proof. It is retained only as a navigation pointer; every load-bearing step is inside the delivered span.\par\medskip
\ifdefined\hypertarget\hypertarget{sourcepointer-6}{}\fi
\noindent\textit{Source reference.} Original lemma 0B6Q: sources/stacks-duality/pinned-original-source.tex, source lines 2171--2171. Same-chapter pointer into the author environment lemma at source line 2171; invoked at proof line 4512 of the counted proof. It is retained only as a navigation pointer; every load-bearing step is inside the delivered span.\par\medskip
\ifdefined\hypertarget\hypertarget{sourcepointer-7}{}\fi
\noindent\textit{Source reference.} Original equation 0A9L: sources/stacks-duality/pinned-original-source.tex, source lines 988--988. Same-chapter pointer into the author environment equation at source line 988; invoked at proof line 4515 of the counted proof. It is retained only as a navigation pointer; every load-bearing step is inside the delivered span.\par\medskip
\ifdefined\hypertarget\hypertarget{sourcepointer-8}{}\fi
\noindent\textit{Source reference.} Original lemma 0B6P: sources/stacks-duality/pinned-original-source.tex, source lines 2119--2119. Same-chapter pointer into the author environment lemma at source line 2119; invoked at proof line 4523 of the counted proof. It is retained only as a navigation pointer; every load-bearing step is inside the delivered span.\par\medskip
\ifdefined\hypertarget\hypertarget{sourcepointer-9}{}\fi
\noindent\textit{Source reference.} Original lemma 0ATX: sources/stacks-duality/pinned-original-source.tex, source lines 4161--4161. Same-chapter pointer into the author environment lemma at source line 4161; invoked at proof line 4553 of the counted proof. It is retained only as a navigation pointer; every load-bearing step is inside the delivered span.\par\medskip
\begin{lemma}
\label{lemma-map-pullback-to-shriek-well-defined}
In Situation \ref{situation-shriek} let
$f : X \to Y$ be a morphism of $\textit{FTS}_S$. There are canonical maps
$$
\mu_{f, K} :
Lf^*K \otimes_{\mathcal{O}_X}^\mathbf{L} f^!\mathcal{O}_Y
\longrightarrow
f^!K
$$
functorial in $K$ in $D^+_\QCoh(\mathcal{O}_Y)$.
If $g : Y \to Z$ is another morphism of $\textit{FTS}_S$, then
the diagram
$$
\xymatrix{
Lf^*(Lg^*K \otimes_{\mathcal{O}_Y}^\mathbf{L} g^!\mathcal{O}_Z)
\otimes_{\mathcal{O}_X}^\mathbf{L} f^!\mathcal{O}_Y
\ar@{=}[d] \ar[r]_-{\mu_f} &
f^!(Lg^*K \otimes_{\mathcal{O}_Y}^\mathbf{L} g^!\mathcal{O}_Z)
\ar[r]_-{f^!\mu_g} &
f^!g^!K \ar@{=}[d] \\
Lf^*Lg^*K \otimes_{\mathcal{O}_X}^\mathbf{L} Lf^* g^!\mathcal{O}_Z
\otimes_{\mathcal{O}_X}^\mathbf{L} f^!\mathcal{O}_Y \ar[r]^-{\mu_f} &
Lf^*Lg^*K \otimes_{\mathcal{O}_X}^\mathbf{L} f^!g^!\mathcal{O}_Z
\ar[r]^-{\mu_{g \circ f}} & f^!g^!K
}
$$
commutes for all $K \in D^+_\QCoh(\mathcal{O}_Z)$.
\end{lemma}

\begin{proof}
If $f$ is proper, then $f^! = a$ and we can use
(\ref{equation-compare-with-pullback}) and if $g$ is also proper,
then Lemma \ref{lemma-transitivity-compare-with-pullback} proves
the commutativity of the diagram (in greater generality).

\medskip\noindent
Let us define the map $\mu_{f, K}$. Choose a compactification
$j : X \to \overline{X}$ of $X$ over $Y$. Since $f^!$ is defined
as $j^* \circ \overline{a}$ we obtain $\mu_{f, K}$ as the restriction
of the map (\ref{equation-compare-with-pullback})
$$
L\overline{f}^*K \otimes_{\mathcal{O}_{\overline{X}}}^\mathbf{L}
\overline{a}(\mathcal{O}_Y)
\longrightarrow
\overline{a}(K)
$$
to $X$. To see this is independent of the choice of the compactification
we argue as in the proof of Lemma \ref{lemma-shriek-well-defined}.
We urge the reader to read the proof of that lemma first.

\medskip\noindent
Assume given a morphism $g : \overline{X}_1 \to \overline{X}_2$
between compactifications $j_i : X \to \overline{X}_i$ over $Y$
such that $g^{-1}(j_2(X)) = j_1(X)$. Denote $\overline{c}$ the
right adjoint for pushforward of Lemma \ref{lemma-twisted-inverse-image}
for the morphism $g$. The maps
$$
L\overline{f}_1^*K \otimes_{\mathcal{O}_{\overline{X}}}^\mathbf{L}
\overline{a}_1(\mathcal{O}_Y)
\longrightarrow
\overline{a}_1(K)
\quad\text{and}\quad
L\overline{f}_2^*K \otimes_{\mathcal{O}_{\overline{X}}}^\mathbf{L}
\overline{a}_2(\mathcal{O}_Y)
\longrightarrow
\overline{a}_2(K)
$$
fit into the commutative diagram
$$
\xymatrix{
Lg^*(L\overline{f}_2^*K \otimes^\mathbf{L}
\overline{a}_2(\mathcal{O}_Y))
\otimes^\mathbf{L} \overline{c}(\mathcal{O}_{\overline{X}_2})
\ar@{=}[d] \ar[r]_-\sigma &
\overline{c}(L\overline{f}_2^*K \otimes^\mathbf{L}
\overline{a}_2(\mathcal{O}_Y)) \ar[r] &
\overline{c}(\overline{a}_2(K)) \ar@{=}[d] \\
L\overline{f}_1^*K \otimes^\mathbf{L} Lg^*\overline{a}_2(\mathcal{O}_Y)
\otimes^\mathbf{L} \overline{c}(\mathcal{O}_{\overline{X}_2})
\ar[r]^-{1 \otimes \tau} &
L\overline{f}_1^*K \otimes^\mathbf{L} \overline{a}_1(\mathcal{O}_Y) \ar[r] &
\overline{a}_1(K)
}
$$
by Lemma \ref{lemma-transitivity-compare-with-pullback}. By
Lemma \ref{lemma-compare-on-open} the maps $\sigma$ and $\tau$
restrict to an isomorphism over $X$. In fact, we can say more.
Recall that in the proof of Lemma \ref{lemma-shriek-well-defined} we used
the map (\ref{equation-sheafy}) $\gamma : j_1^* \circ \overline{c} \to j_2^*$
to construct our isomorphism
$\alpha_g : j_1^* \circ \overline{a}_1 \to j_2^* \circ \overline{a}_2$.
Pulling back to map $\sigma$ by $j_1$ we obtain the identity
map on $j_2^*\left(L\overline{f}_2^*K \otimes^\mathbf{L}
\overline{a}_2(\mathcal{O}_Y)\right)$ if we identify
$j_1^*\overline{c}(\mathcal{O}_{\overline{X}_2})$
with $\mathcal{O}_X$ via $j_1^* \circ \overline{c} \to j_2^*$, see
Lemma \ref{lemma-restriction-compare-with-pullback}.
Similarly, the map $\tau : Lg^*\overline{a}_2(\mathcal{O}_Y)
\otimes^\mathbf{L} \overline{c}(\mathcal{O}_{\overline{X}_2}) \to
\overline{a}_1(\mathcal{O}_Y) = \overline{c}(\overline{a}_2(\mathcal{O}_Y))$
pulls back to the identity map on $j_2^*\overline{a}_2(\mathcal{O}_Y)$.
We conclude that pulling back by $j_1$ and applying $\gamma$ wherever
we can we obtain a commutative diagram
$$
\xymatrix{
j_2^*\left(L\overline{f}_2^*K \otimes^\mathbf{L}
\overline{a}_2(\mathcal{O}_Y)\right) \ar[r] \ar[d] &
j_2^*\overline{a}_2(K) \\
j_1^*L\overline{f}_1^*K \otimes^\mathbf{L} j_2^*\overline{a}_2(\mathcal{O}_Y) &
j_1^*(L\overline{f}_1^*K \otimes^\mathbf{L} \overline{a}_1(\mathcal{O}_Y))
\ar[r] \ar[l]_{1 \otimes \alpha_g} &
j_1^* \overline{a}_1(K) \ar[lu]_{\alpha_g}
}
$$
The commutativity of this diagram exactly tells us that the map
$\mu_{f, K}$ constructed using the compactification $\overline{X}_1$
is the same as the map $\mu_{f, K}$ constructed using the compactification
$\overline{X}_2$ via the identification $\alpha_g$ used in the proof
of Lemma \ref{lemma-shriek-well-defined}. Some categorical arguments
exactly as in the proof of  Lemma \ref{lemma-shriek-well-defined}
now show that $\mu_{f, K}$ is well defined (small detail omitted).

\medskip\noindent
Having said this, the commutativity of the diagram in the statement
of our lemma follows from the construction of the isomorphism
$(g \circ f)^! \to f^! \circ g^!$ (first part of the proof of
Lemma \ref{lemma-upper-shriek-composition} using
$\overline{X} \to \overline{Y} \to Z$) and the result
of Lemma \ref{lemma-transitivity-compare-with-pullback}
for $\overline{X} \to \overline{Y} \to Z$.
\end{proof}

\end{document}
